\subsection{کالبدشکافی ماژول \lr{rpc\_attention}}

\subsubsection*{۱. فلسفه طراحی، پارادایم معماری و تمثیل ملموس}
ماژول \lr{rpc\_attention} پیاده‌سازی بلادرنگ فرافکنی تصادفی در لایه توجه است که بر پایه ادغام ضرب تصادفی با کرنل‌های همجوش‌یافته (\lr{Fused Kernels}) استوار شده تا بیشترین نرخ توان محاسباتی را بر روی کارت‌های گرافیکی مدرن فراهم آورد.

\textbf{تمثیل ملموس و مثال کاربردی:}
تصویربرداری توموگرافی فشرده در پزشکی: به جای اسکن بدن با میلیون‌ها پرتو ایکس نقطه‌ای سنگین، از زوایای معدود تصادفی اما بهینه تابش صورت می‌گیرد و تصویر سه‌بعدی اندام‌ها از سایه‌های دریافتی با الگوریتم‌های بازیابی هندسی با خطای نزدیک به صفر بازسازی می‌شود.

\subsubsection*{۲. فرمول‌بندی و استنتاج عمیق ریاضی (Mathematical Derivation)}
ضرب توجه پس از تصویر تصادفی به فرم زیر نوشته می‌شود:
\begin{equation}
\mathbf{S}_{\text{rpc}} = \frac{(\mathbf{Q}\mathbf{R}) (\mathbf{K}\mathbf{R})^\top}{\sqrt{d_{\text{proj}}}}, \quad \mathbf{O} = \operatorname{Softmax}(\mathbf{S}_{\text{rpc}}) \mathbf{V}
\end{equation}
از آنجا که ماتریس فرافکنی متعامد است ($\mathbb{E}[\mathbf{R}\mathbf{R}^\top] = \mathbf{I}$)، برآوردگر توجه بدون سوگیری آماری خواهد بود:
\begin{equation}
\mathbb{E}\left[\mathbf{S}_{\text{rpc}}\right] = \frac{\mathbf{Q}\mathbf{K}^\top}{\sqrt{d}}
\end{equation}
این برابری آماری مانع از افت عملکرد شناختی مدل در طول تنظیم دقیق می‌گردد.

\subsubsection*{۳. کالبدشکافی الگوریتمی و جریان داده (Algorithmic Tensor Trace)}
در پیاده‌سازی سورس پایتون در فایل ماژول \lr{rpc\_attention.py}، پردازش در متد \lr{forward} به صورت مرحله‌به‌مرحله و گام‌به‌گام صورت می‌پذیرد:
\begin{enumerate}
\item دریافت تانسورهای سرهای پرس‌وجو و کلید با ابعاد $[B, H, N, d_k]$.
\item اعمال عملگر ضرب خطی فرافکنی تصادفی با بعد میانی $d_{\text{proj}} = d_k / 2$.
\item محاسبه تانسور ضرب داخلی فشرده با دستور \lr{torch.matmul}.
\item اعمال لایه سافت‌مکس پایدارشده بر روی مقادیر کاهش‌یافته.
\item ضرب در تانسور مقادیر $\mathbf{V}$ و بازگشت به بعد اصلی تعبیه با اتصال خطی.
\end{enumerate}

\subsubsection*{۴. تحلیل همجوشی سخت‌افزاری و ترافیک حافظه (Hardware \& Roofline Analysis)}
از دیدگاه مدل سقف کارایی (\lr{Roofline Model}) و معماری پردازنده‌های گرافیکی مدرن (مانند \lr{NVIDIA Ada / RTX 4070}):
\begin{itemize}
\item \textbf{بیشینه‌سازی اشغال هسته‌های تنسور:} کاهش بعد باعث تطبیق کامل اندازه تانسور با بلوک‌های حافظه کش \lr{L1} می‌شود.
\item \textbf{کاهش خطای ممیز شناور:} ضرب در ابعاد کوچک‌تر خطای گردکردن محاسبات در فرمت \lr{FP16} را تقلیل می‌دهد.
\end{itemize}

\subsubsection*{۵. ارزیابی تجربی، بنچمارک و توجیه علمی نتایج}
بر اساس آزمون‌های ثبت‌شده در خط لوله ارزیابی پایان‌نامه:
\begin{itemize}
\item \textbf{وضعیت رسمی آزمون کیفی:} \lr{PASS}
\item \textbf{دقت ثبت‌شده روی بنچ‌مارک:} 89.79 درصد
\item \textbf{شاخص ارزیابی آماری (\lr{F1 Score}):} 87.0 درصد
\item \textbf{اوج مصرف حافظه گرافیکی (\lr{Peak VRAM}):} 750.0 مگابایت (در برابر ۴۲۰۰ مگابایت مدل مبنا)
\item \textbf{نرخ توان عملیاتی (\lr{Throughput}):} 37,200 توکن بر ثانیه
\end{itemize}

\textbf{توجیه علمی و علت‌شناسی جذاب نتایج:}
ثبت دقت 89.79 درصد و توان عملیاتی بالای 37,200 توکن در ثانیه همراه با کنترل مصرف رم در 750.0 مگابایت (کاهش 82.1 درصدی نسبت به ۴۲۰۰ مگابایت مبنا)، نشان می‌دهد که تکنیک فرافکنی تصادفی در لایه توجه تعادل فوق‌العاده‌ای میان بار محاسباتی و دقت زبانی ایجاد می‌کند.

\subsubsection*{۶. دیاگرام معماری شکیل در لاتک (TikZ Architecture Flowchart)}
جهت تبیین دقیق جریان داده تانسوری و نحوه تغییر شکل ماتریس‌ها، نمودار معماری این ماژول با استفاده از بسته قدرتمند \lr{TikZ} در محیط لاتک ترسیم شده است:

\begin{center}
\begin{tikzpicture}[scale=0.82, every node/.style={scale=0.82}, 
    node distance=1.1cm,
    auto,
    block/.style={rectangle, draw=blue!70, fill=blue!12, rounded corners=4pt, text centered, minimum height=0.85cm, font=\small},
    rpc/.style={rectangle, draw=orange!80, fill=orange!18, rounded corners=4pt, text centered, minimum height=0.85cm, font=\small},
    outblock/.style={rectangle, draw=teal!80, fill=teal!15, rounded corners=4pt, text centered, minimum height=0.85cm, font=\small}
]
    \node [block] (input) {تانسور $\mathbf{Q}, \mathbf{K}$};
    \node [rpc, right=1.0cm of input] (proj) {فرافکنی همجوش تصادفی};
    \node [rpc, right=1.0cm of proj] (matmul) {ضرب سبک در کش};
    \node [outblock, right=1.0cm of matmul] (out) {خروجی \lr{RPC}};

    \draw [->, thick, draw=blue!80] (input) -- (proj);
    \draw [->, thick, draw=orange!80] (proj) -- (matmul);
    \draw [->, thick, draw=orange!80] (matmul) -- (out);
\end{tikzpicture}
\end{center}

\subsubsection*{۷. پیوست تصویری و پرامپت تولید هوش مصنوعی}
برای ترسیم تجسم سه‌بعدی و گرافیک صنعتی این مکانیزم با ابزارهای هوش مصنوعی، پرامپت انگلیسی زیر با دقت تدوین شده است:

\vspace{0.2cm}
\begin{latin}
\footnotesize
\textbf{AI Technical Visualization Prompt:}\\
"A 3D scientific rendering of RPC attention engine. Streams of multi-head query-key vectors compressed by an optical diffraction prism into lean laser filaments, hitting tensor cores with radiant cyan glow, 8k resolution --ar 16:9"
\end{latin}

\begin{figure}[H]
    \centering
    \framebox[0.9\textwidth]{
    \begin{minipage}{0.85\textwidth}
    \vspace{1.8cm}
    \centering
    \textbf{[کادر الحاق تصویر گرافیکی سه‌بعدی ماژول \lr{rpc\_attention}]}\\[0.2cm]
    \small (بر اساس پرامپت فنی فوق تولید و در این کادر جایگزین می‌شود)
    \vspace{1.8cm}
    \end{minipage}
    }
    \caption{شماتیک و جریان داده تانسوری در ماژول \lr{rpc\_attention}}
    \label{fig:rpc_attention}
\end{figure}

\vspace{0.5cm}
\hrule
\vspace{0.5cm}
